\clearpage
\SourceLeaf{154}\pdfbookmark[1]{Abschnitt IV.}{sectio-quarta}
\begin{center}
{\large ABSCHNITT IV.}\par\medskip
\emph{Beziehungen zwischen mehreren Orten im Raum.}\par\bigskip
\hypertarget{native-154-article110}{}\pdfbookmark[2]{110.}{article-110}110.
\end{center}
Die in diesem Abschnitt betrachteten Beziehungen sind von der Beschaffenheit der Bahn unabhängig und beruhen allein auf der Annahme, dass alle Bahnpunkte mit der Sonne in derselben Ebene liegen. Wir begnügen uns hier damit, einige der einfachsten Beziehungen zu berühren, und behalten uns die verwickelteren und spezielleren für das zweite Buch vor.

Zwei Orte eines Himmelskörpers im Raum bestimmen die Lage der Bahnebene vollständig, sofern sie nicht mit der Sonne auf derselben Geraden liegen. Da sich der Ort eines Punktes im Raum vor allem auf zwei Arten angeben lässt, ergeben sich daraus zwei Aufgaben.

Wir nehmen zunächst an, die beiden Orte seien durch ihre heliozentrischen Längen und Breiten gegeben, die wir in dieser Reihenfolge mit $\lambda$, $\lambda'$, $\beta$ und $\beta'$ bezeichnen. Die Abstände von der Sonne gehen in die Rechnung nicht ein. Bezeichnet $\AscendingNode$ die Länge des aufsteigenden Knotens und $i$ die Neigung der Bahn gegen die Ekliptik, so gilt
\[
\begin{aligned}
\operatorname{tang}\beta&=\operatorname{tang}i\sin(\lambda-\AscendingNode)\\
\operatorname{tang}\beta'&=\operatorname{tang}i\sin(\lambda'-\AscendingNode).
\end{aligned}
\]
Die Bestimmung der Unbekannten $\AscendingNode$ und $\operatorname{tang}i$ führt hier auf die in Artikel 78, II behandelte Aufgabe zurück. Nach der ersten Lösung erhalten wir
\[
\begin{aligned}
\operatorname{tang}i\sin(\lambda-\AscendingNode)&=\operatorname{tang}\beta\\
\operatorname{tang}i\cos(\lambda-\AscendingNode)&=\frac{\operatorname{tang}\beta'-\operatorname{tang}\beta\cos(\lambda'-\lambda)}{\sin(\lambda'-\lambda)};
\end{aligned}
\]
\SourceLeaf{155}Nach der dritten Lösung hingegen finden wir $\AscendingNode$ aus der Gleichung
\[
\operatorname{tang}(\tfrac12\lambda+\tfrac12\lambda'-\underset{\rule{.45em}{.35pt}}{\AscendingNode})=\frac{\sin(\beta'+\beta)\operatorname{tang}\tfrac12(\lambda'-\lambda)}{\sin(\beta'-\beta)}
\]
Das ist etwas bequemer, wenn die Winkel $\beta$ und $\beta'$ unmittelbar gegeben sind und nicht durch die Logarithmen ihrer Tangenten. Zur Bestimmung von $i$ müssen wir jedoch auf eine der folgenden Formeln zurückgreifen
\[
\operatorname{tang}i=\frac{\operatorname{tang}\beta}{\sin(\lambda-\AscendingNode)}=\frac{\operatorname{tang}\beta'}{\sin(\lambda'-\AscendingNode)}\mathpunct{\cdot}
\]
Bei der Bestimmung des Winkels $\lambda-\AscendingNode$ oder $\tfrac12\lambda+\tfrac12\lambda'-\AscendingNode$ aus seiner Tangente ist die Mehrdeutigkeit so aufzulösen, dass $\operatorname{tang}i$ positiv oder negativ ausfällt, je nachdem, ob die auf die Ekliptik projizierte Bewegung rechtläufig oder rückläufig ist. Diese Unsicherheit lässt sich also nur beheben, wenn feststeht, auf welcher Seite der Himmelskörper vom ersten Ort zum zweiten gelangt ist. Andernfalls kann man den aufsteigenden Knoten nicht vom absteigenden unterscheiden.

Sind die Winkel $\AscendingNode$ und $i$ gefunden, so erhält man die Breitenargumente $u'$ und $u$ aus den Formeln
\[
\operatorname{tang}u=\frac{\operatorname{tang}(\lambda-\AscendingNode)}{\cos i},\qquad
\operatorname{tang}u'=\frac{\operatorname{tang}(\lambda'-\AscendingNode)}{\cos i},
\]
Die Breitenargumente sind im ersten oder zweiten Halbkreis zu nehmen, je nachdem, ob die entsprechenden Breiten nördlich oder südlich sind. Diesen Formeln fügen wir noch die folgenden hinzu, von denen man nach Belieben die eine oder andere zur Kontrolle der Rechnung verwenden kann.
\[
\begin{aligned}
\cos u&=\cos\beta\cos(\lambda-\AscendingNode),&\quad \cos u'&=\cos\beta'\cos(\lambda'-\AscendingNode),\\[3pt]
\sin u&=\frac{\sin\beta}{\sin i},&\sin u'&=\frac{\sin\beta'}{\sin i},
\end{aligned}
\]
\[
\begin{aligned}
\sin(u'+u)&=\frac{\sin(\lambda+\lambda'-2\AscendingNode)\cos\beta\cos\beta'}{\cos i},\\[3pt]
\sin(u'-u)&=\frac{\sin(\lambda'-\lambda)\cos\beta\cos\beta'}{\cos i}\mathpunct{\cdot}
\end{aligned}
\]
